
\bigskip

\section{The heq construction and the exact completion}\label{heq}
We have now established in Theorem \ref{eliminate hyper iff pro exact} a correspondence between two properties---elimination of hyperimaginaries and the exactness of $\Pro(\defT(T)_{ex/reg})$. Just as Theorem \ref{Teq is exact} is a constructive version of Proposition \ref{EI iff exact}, in this section we show there is a constructive version of Theorem \ref{eliminate hyper iff pro exact}. Indeed, we will show that the heq construction corresponds to the ex/reg completion of $\Pro(\defT(T))$.

\medskip

\subsection{The heq construction}
We recall the heq construction from \cite{DobrowolskiJan2022Kipl}. 

As notation, given a type-definable equivalence relation $E(x,x')$, we denote by $[x]_E$ the $E$-hyperimaginary corresponding to the $E$-class of $x$. When it is clear, we omit the subscript and write $[x]$ for $[x]_E$. If $y$ is variable of the sort of an $E$-hyperimaginary, we denote by $y_r$ a variable of the home sort corresponding to a representative. Standard model theoretic notation dictates that one can write $a\in\M$ when $a$ is a finite tuple in $\M$; we shall extend this abuse of notation to the case when $a$ is a small infinite tuple.



\begin{defin}[{\cite[Definitions 10.7 and 10.8]{DobrowolskiJan2022Kipl}}]
    Let $\M$ be a model of $T$. Let $\E$ be a set of type-definable equivalence relations.
    \begin{enumerate}
        \item The language $\Lang_\E$ expands the original language $\Lang$ by sorts $S_E$ for every $E\in\E$. For every $\Lang$-formula $\phi(x,y)$ and type-definable equivalence relatio $E$ on the sort of $y$, we add in a new predicate $R_\phi(x,[y]_E)$. However, we consider $\Lang_\E$ as a language in positive logic, even though $T$ is Boolean. 
        \item The $\Lang_\E$-structure $\M^\E$ expands $\M$ as follows: the sort $S_E$ is interpreted as the quotient of $E(\M)$. For each $\Lang$-formula $\phi(x,y)$, type-definable equivalence relation $E$ on the sort of $y$, and elements $a,b\in\M$,
        \[ \M^\E \models R_\phi(a,[b]_E) \ \Leftrightarrow \ \textup{there is }b' \textup{ such that } \M\models \phi(a,b')\wedge E(b,b'). \]
    \end{enumerate}
\end{defin}

From now on, fix a monster model $\mathcal{U}$ of $T$.

One would like to add in all the hyperimaginaries, but this is a proper class since the variables can be arbitrarily long. Luckily, we have the following lemma. 

\begin{lem}[{\cites[Corollary 3.3\addsemicolon]{Ben-Yaacov2003c}[Fact 1.1]{Hart2000-HARCAC-24}}]
    Let $E(x,y)$ be a type-definable equivalence relation. Then there is a set $I$ and type-definable equivalence relations $E_i(x_i,y_i) : i\in I$, where $x_i$ is a subtuple of $x$ with $|x_i|\leq |T|$, such that for all $a,b\in\mathcal{U}$,
    \[ \mathcal{U}\models E(a,b) \ \Leftrightarrow \ \mathcal{U}\models E_i(a_i,b_i) \textup{ for all }i\in I. \]
    Moreover, if $T$ is Boolean, we may assume $|x_i|$ is countable.
\end{lem}

\begin{defin}
    We define $heq$ to be the set of type-definable equivalence relations $E(x,y)$, where $|x|\leq|T|$.
\end{defin}


The following is the key lemma that allows us to describe $\mathcal{U}^{heq}$ using $\mathcal{U}$.

\begin{lem}[{\cite[Lemma 10.10]{DobrowolskiJan2022Kipl}}]\label{heq translate}
    Let $\phi(x)$ be an $\Lang_{heq}$-formula. Then there is a set $\Sigma(x_r)$ of $\Lang$-formulas such that for all $a\in\mathcal{U}^{heq}$ and $a_r\in\mathcal{U}$ representing $a$,
    \[ \mathcal{U}^{heq}\models\phi(a) \ \Leftrightarrow \ \mathcal{U}\models\Sigma(a_r). \]
\end{lem}

\begin{remark}[{\cite[Lemma 10.11]{DobrowolskiJan2022Kipl}}]
    It is clear that the lemma works even when we replace the $\Lang_{heq}$-formula with a set of $\Lang_{heq}$-formulas.
\end{remark}

Shelah's original eq construction specifically expands a structure with imaginaries, so that this expansion eliminates imaginaries. This is not the case with heq and hyperimaginaries. One does not apply the heq construction, so that the expanded structure eliminates hyperimaginaries. Rather, the heq construction simply expands a structure \emph{with} hyperimaginares.

\begin{remark}[{\cite[Lemma 10.13]{DobrowolskiJan2022Kipl}}]\label{typedef quotient heq}
    Let $E(x_r,y_r)$ be a type-definable equivalence relation in $T$ with $|x_r|=I$. Then, the quotient map $\mathcal{U}^I \rightarrow S_E$ is type-definable in $\mathcal{U}^{heq}$. This is given as follows: for every formula $\phi(x_r,y_r) \in E(x_r,y_r)$, we obtain a predicate $R_\phi(x_r,y)$. We claim that the quotient map is precisely the realisations of the type $P:= \{R_\phi(x_r,y)\ |\ \phi\in E\}$.

    To see this, assume $\mathcal{U}^{heq}\models R_\phi(a,[a'])$ for all $\phi\in E$, for some $a,a'\in\mathcal{U}^I$. We wish to show that $[a']$ is the representative of $a$. Consider the type
    \[ \Gamma(y_r) = E(a,y_r) \cup E(y_r,a'). \]
    We show $\Gamma$ is finitely satisfiable. Given $\phi\in E$, we have $\mathcal{U}^{heq}\models R_\phi(a,[a'])$, so there is some $a''\in\mathcal{U}$ such that $\mathcal{U}\models \phi(a,a'')\wedge E(a'',a')$. By compactness and saturation, there is $b\in\mathcal{U}$ realising $\Gamma$. Then, $[a]=[b]=[a']$ as required. This shows that the realisations of $P$ are contained in the quotient map.

    For the converse, assume $\mathcal{U}\models E(a,a')$. In particular, $\mathcal{U}^{heq}\models R_\phi(a,[a'])$ for all $\phi\in E$.
\end{remark}

It was shown in \cite[Example 3.1.6]{alma992983361725701631} that $\textup{Th}(\mathcal{U}^{heq})$ does not axiomatise the class $\{ \N^{heq} \ | \ \N\models T\}$. In particular, the type $P$ in Remark \ref{typedef quotient heq} might not define a function (let alone the quotient map) in some model $\N\models\textup{Th}(\mathcal{U}^{heq})$. Otherwise, $P$ would be a morphism in $\Pro(\defT(\textup{Th}(\mathcal{U}^{heq})))$ and thus a cofiltered limit of functions definable in $\textup{Th}(\mathcal{U}^{heq})$. In other words, we would require there to be a subset $E_0\subseteq E$ such that $\{R_\phi\ | \ \phi\in E_0\}$ defines the same set as $P$ and for every $\phi\in E_0$, the predicate $R_\phi$ defines a function. The following example shows this does not happen.

\begin{example}
    Let $T$ be the theory of infinite sets in the empty language. Let $E$ be the discrete equivalence relation on tuples $(x_n)_{n\in\omega}$: 
    \[E((x_n),(y_n)) \ \Leftrightarrow \ x_n=y_n \textup{ for all }n.\]
    Then, the sort $S_E$ is interpreted in $\U^{heq}$ by $\U^\omega$. A typical formula $\phi(x_1,\ldots,x_m,(y_n)_{n\in\omega})\in E$ has the form $\bigwedge_{i=1}^m x_i=y_i$. The interpretation of $R_\phi(x_1,\ldots,x_m,y)$ in $\U^{heq}$ is that $(x_1,\ldots,x_m)$ is a subtuple of $y$ (which is an element of $\U^\omega$). In particular, given $a_1,\ldots,a_m$, there are infinitely many $b$ such that $R_\phi(a_1,\ldots,a_m,b)$ holds, so $R_\phi$ does not define a function. Indeed, if we fix a sequence $(a_n)_{n\in\omega}$, the type $\{R_\phi(a_1,\ldots,a_m,y)\wedge R_\phi(a_1,\ldots,a_m,z)\wedge y\neq z \ | \ \phi\in E\}$ is consistent.
\end{example}

When we first learn about non-standard models of arithmetic, we might find the existence of infinite elements and its relative consistency with the axioms of arithmetic puzzling. Na\"{i}vely, we might try to axiomatise away these non-standard models. Perhaps we expand the structure by a predicate $U$, whose interpretation is all the standard integers, and extend our theory by the sentence saying that all elements are $U$. As we know, this does not work: non-standard models always exist. Passing from $\U$ to $\U^{heq}$ is exactly doing this. In hindsight, it is no surprise that $\textup{Th}(\U^{heq})$ does not axiomatise that the type $P$ from Remark~\ref{typedef quotient heq} defines a function. This is related to the fact that $\Pro(\Pro(\C))$ is in general not equivalent to $\Pro(\C)$. Likewise, $(\U^{heq})^{heq}$ is not the same as $\U^{heq}$ in any sense. 

If we wish to build a syntactic category of definable sets and functions from the structure $\U^{heq}$, we would want types like $P$ above to be morphisms. We have seen that $P$ is not a morphism in $\Pro(\defT(\textup{Th}(\U^{heq})))$. The problem lies in the fact that when we take the theory $\textup{Th}(\U^{heq})$, we lose information which is type-definable, for example that $P$ defines a function in $\U^{heq}$. Hence, when we extract a theory from $\U^{heq}$, we should instead do so in an infinitary positive logic. Consider that $\mathcal{U}^{heq}$ essentially is the structure of hyperdefinable sets of $\mathcal{U}$, where a set is hyperdefinable if it is the quotient of some type-definable equivalence relation. These hyperdefinable sets are closed under infinite conjunctions, infinite existential quantification and finite disjunction, but \emph{not} under negation. This corresponds exactly to infinitary positive logic, or the internal logic of a $\infty$-coherent category. We pursue this in the next subsection.

\medskip

\subsection{The infinitary syntactic category of a structure}
We study an infinitary positive variant of $\defT(T)$. We first define this semantically on a model.

\begin{defin}
    We define the \emph{infinitary syntactic category} $\tpdefT(\M)$ of a structure $\M$ as follows:
    \begin{itemize}
        \item Objects are type-definable sets of $\M$;
        \item Morphisms are type-definable sets of $\M$, which are the graphs of some function. 
    \end{itemize}
\end{defin}

Two different models of the same theory can give rise to different infinitary syntactic categories, but if the two models are sufficiently saturated, then the infinitary syntactic categories are equivalent. 

\begin{prop}
    Let $\M$ and $\N$ be monster models of the same theory $T$. Let $\tpdefT(\M)$ and $\tpdefT(\N)$ denote their infinitary syntactic categories. Then $\tpdefT(\M) \simeq \tpdefT(\N)$. 
\end{prop}

\begin{proof}
    The required functor sends the realisation of a type $X$ in $\M$ to the realisation of $X$ in $\N$. To ensure this is well defined, we prove that if $X$ and $Y$ are types with the same realisations in $\M$, then they have the same realisations in $\N$.

    Suppose not, say there is some $a\in\N$ realising $X$ but not $Y$. Then, there is some formula $\psi(x)$ satisfied by $a$, such that $q\vdash\neg\psi$. Then, the type $p\cup\{\psi\}$ is consistent and realised by some $b$ in $\M$. Then, $b$ realises $X$ but not $Y$ in $\M$.

    Next, we prove that if $P$ is a type-definable function from the type $X$ to the type $Y$ in $\M$, then this is also the case in $\N$. Suppose this is not the case in $\N$.

    \textit{Case} 1:\quad In $\N$, the set defined by $P$ is not total. Consider that the types $X$ and $\exists y \ (Y(y)\wedge X(x)\wedge P(x,y))$ are equivalent in $\M$. By above, they are also equivalent in $\N$. Hence, every $a\in X(\N)$ has an image in $Y$ under $P$.

    \textit{Case} 2:\quad In $\N$, the set defined by $P$ is not single-valued. Let $a\in X(\N)$ be such that there are distinct $b,b'\in Y(\N)$ and $P(a,b)\wedge P(a,b)$ holds. Let $y_0$ be the coordinate where $b,b'$ differ. Consider that the type $X(x)\wedge Y(y)\wedge Y(y')\wedge y_0\neq y'_0\wedge P(x,y)\wedge P(x,y')$ is consistent and realised in $\M$. Hence, $P$ is not single-valued in $\M$ either.
\end{proof}

Given that $\U$ is sufficiently saturated, the two categories $\tpdefT(\mathcal{U})$ and $\Pro(\defT(\textup{Th}(\mathcal{U})))$ are equivalent. This is an easy corollary of Lemma~\ref{functions of Pro}.

\begin{example}\label{positive not work ex}
In the Boolean setting, we can always present a morphism of $\tpdefT(\mathcal{U})$ as a morphism of $\Pro(\defT(T))$ by using compactness of non-positive formulas. This is in contrast to the setting when $T$ is a theory in positive logic. Let $T$ be the empty positive theory in the empty language with a single sort $S$. It is known that the positively closed models of $T$ are precisely the singleton sets \cite[Remark~2.1.12]{kamsma2025positivelogicintroductionmodel}. Hence, the formula $x=x$ is interpreted as the graph of a function $\top\rightarrow S$. However, $T$ does not prove that it is a function.
\end{example}

\medskip

\subsection{Connection with the exact completion}
We return to the case of $\mathcal{U}^{heq}$. 

\begin{thm}
    The category $\tpdefT(\U^{heq})$ is $\infty$-exact. 
\end{thm}

\begin{proof}
    The proof that $\tpdefT(\U^{heq})$ is $\infty$-regular is routine.

    We saw in Remark \ref{typedef quotient heq} that if an equivalence relation is type-definable in $T$, then it has a type-definable quotient in $\tpdefT(\U^{heq})$. 

    Now let $R([x_r]_E,[y_r]_E)$ be a type-definable equivalence relation in $\mathcal{U}^{heq}$, where $E$ is some type-definable equivalence relation in $T$. By Lemma \ref{heq translate}, there is an equivalence relation $\Gamma(x_r,y_r)$ type-definable in $T$, such that for all $a_r,b_r\in\mathcal{U}$,
    \[ \mathcal{U}\models\Gamma(a_r,b_r) \ \Leftrightarrow \ \mathcal{U}^{heq}\models R([a_r]_E,[b_r]_E). \]

    Now $\Gamma$ has a corresponding sort $S_\Gamma$ in $\mathcal{U}^{heq}$. Moreover, this comes with a type-definable quotient map $g: \mathcal{U}^I\rightarrow S_\Gamma$, where $I = |x_r|$. Similarly, there is a type-definable quotient map $f:\mathcal{U}^I\rightarrow S_E$. We define the quotient map $h:S_E \rightarrow S_\Gamma$ by
    \[ h([x_r]_E) = [y_r]_\Gamma \ \Leftrightarrow \ \exists z_r \ (f(z_r)=[x_r]_E \wedge g(z_r) = [y_r]_\Gamma). \]
    Type-definable sets are closed under infinite conjunection and existential quantification, so $h$ is type-definable indeed.
\end{proof}

Since every $T$-definable set is type-definable in $\U^{heq}$, there is a natural inclusion $\defT(T) \hookrightarrow \tpdefT(\U^{heq})$. Moreover, by compactness, this inclusion is full. By the universal property of the pro-completion and the exact completion, we obtain the following induced $\infty$-regular functors.

% https://q.uiver.app/#q=WzAsNCxbMCwwLCJcXGRlZlQoVCkiXSxbMSwwLCJcXGRlZlReXFxpbmZ0eShUKSJdLFsyLDAsIihcXGRlZlReXFxpbmZ0eShUKSlfe2V4L3JlZ30iXSxbMiwxLCJcXEMiXSxbMCwxLCIiLDIseyJzdHlsZSI6eyJ0YWlsIjp7Im5hbWUiOiJob29rIiwic2lkZSI6InRvcCJ9fX1dLFsxLDIsIiIsMix7InN0eWxlIjp7InRhaWwiOnsibmFtZSI6Imhvb2siLCJzaWRlIjoidG9wIn19fV0sWzAsMywiIiwwLHsiY3VydmUiOjJ9XSxbMSwzLCIiLDEseyJzdHlsZSI6eyJib2R5Ijp7Im5hbWUiOiJkYXNoZWQifX19XSxbMiwzLCJGIiwwLHsic3R5bGUiOnsiYm9keSI6eyJuYW1lIjoiZGFzaGVkIn19fV1d
\[\begin{tikzcd}
	{\defT(T)} & {\Pro(\defT(T))} & {(\Pro(\defT(T)))_{ex/reg}} \\
	&& \tpdefT(\U^{heq})
	\arrow[hook, from=1-1, to=1-2]
	\arrow[curve={height=12pt}, from=1-1, to=2-3]
	\arrow[hook, from=1-2, to=1-3]
	\arrow[dashed, from=1-2, to=2-3]
	\arrow["F", dashed, from=1-3, to=2-3]
\end{tikzcd}\]

For notation, we consider the functors $\defT(T)\hookrightarrow\Pro(\defT(T))\hookrightarrow(\Pro(\defT(T)))_{ex/reg}$ as inclusions. One can think of the functor $F$ as follows. An object of $\defT(T)$, which is an $\Lang$-formula, is sent to its realisations in $\mathcal{U}^{heq}$. An object of $\Pro(\defT(T))$, which is a type of $T$, is sent to the cofiltered limit of the realisations of its formulas. From \cite[Corollary 10]{KAMENSKY2007180}, we know this is precisely the realisations of the type itself. An object of $(\Pro(\defT(T)))_{ex/reg}$, which is a quotient of a type by a type-definable equivalence relation in $T$, is sent to its corresponding quotient in $\tpdefT(\U^{heq})$. 

\begin{thm}
    The functor $F:(\Pro(\defT(T)))_{ex/reg} \rightarrow \tpdefT(\U^{heq})$ is an equivalence of categories.
\end{thm}

\begin{proof}
    We prove that $F$ is essentially surjective and fully faithful.
    
    Essential surjectivity: \quad Let $\Gamma([x]_E)$ be a type in $T^{heq}$, where $E$ is an equivalence relation on $\mathcal{U}^I$ type-definable in $T$. Then there is a type $\Sigma(x)$ in $T$ such that for all $a\in\mathcal{U}$
    \[ \Gamma([a]_E) \ \Leftrightarrow \ \Sigma(a). \]
    Define $E\upharpoonright_\Sigma$ to be the equivalence relation which is $E$ restricted to $\Sigma$. Categorically, this is the pullback $E \times_{\mathcal{U}^I\times\mathcal{U}^I} (\Sigma\times\Sigma)$. It is not hard to see that the following is a quotient diagram. 
    \[F(E\upharpoonright_\Sigma) \rightrightarrows F(\Sigma) \rightarrow \Gamma\]
    In $(\Pro(\defT(T)))_{ex/reg}$, there is $\textup{coeq}(E\upharpoonright_\Sigma \rightrightarrows \Sigma)$. Now, $F$ preserves quotients, so we have
    \[F(\textup{coeq}(E\upharpoonright_\Sigma \rightrightarrows \Sigma)) \cong \Gamma.\]

\smallskip

    Fullness: \quad Let $f\in\Hom_{\tpdefT(\U^{heq})}(F(X/D),F(Y/E))$, where $X/D,Y/E$ denote the quotients of $X,Y$ by equivalence relations $D,E$ respectively, and $X,Y,D,E\in\Pro(\defT(T))$. In other words, $f$ is type-definable in $\mathcal{U}^{heq}$. By Lemma \ref{heq translate}, there is a relation $R$ from $X$ to $Y$ which is type-definable in $T$ and such that for all $a,b\in\mathcal{U}$,
    \begin{equation}
        R(a,b) \ \Leftrightarrow \ f([a]_D)=[b]_E. \tag{$\dagger$}
    \end{equation}
    Now, we can consider $R$ as an object in $\Pro(\defT(T))$ and check that it satisfies Definition \ref{ex/reg carboni}. 
    
    First, we show that $R = RD$. Note that $RD(x,y)$ is the type $\exists x' \ D(x,x') \wedge R(x',y)$. By Proposition \ref{isomorphism is calculated in monster}, it is sufficient to show that $RD$ and $R$ have the same realisations in $\mathcal{U}$. That $R(\mathcal{U})\subseteq RD(\mathcal{U})$ follows from the reflexivity of $R$. Conversely, if $D(a,a')\wedge R(a',b)$ holds for some $a,a',b\in\mathcal{U}$, then $f([a]_D) = f([a']_D) = [b]_E$. Thus, we have $R(a,b)$. Similarly, one may show that $R = ER$.

    Second, we show that $D\leq R^\circ R$. Now, $R^\circ R(x,x')$ is the type $\exists y \ (R(x,y)\wedge R(x',y))$. It suffices to show that there is a monomorphism $D(x,x') \hookrightarrow R^\circ R(x,x')$ in $\Pro(\defT(T))$ given $(x,x')\mapsto (x,x')$. Now, this is a type-definable set, whose realisations form the graph of a function from $D(\mathcal{U})$ to $R^\circ R(\mathcal{U})$. By Lemma \ref{functions of Pro}, this corresponds to a morphism $\iota\in\Hom_{\Pro(\defT(T))}(D,R^\circ R)$. 
    To show that $\iota$ is a monomorphism, we consider $\Pro(\defT(T))$ as a subcategory of \textbf{Set} by interpreting in $\mathcal{U}$. Certainly, $\iota(\mathcal{U})$ is injective, so is a monomorphism in \textbf{Set}. Hence, it is also a monomorphism in $\Pro(\defT(T))$. Similarly, one may show that $RR^\circ \leq E$.
    
    Hence, $R\in\Hom_{(\Pro(\defT(T)))_{ex/reg}}(X/D,Y/E)$. For notation, we write $\tilde{f}$ for this morphism and reserve $R$ for the object.
    We claim that $F(\tilde{f})=f$. The condition $(\dagger)$ means that in ${\tpdefT(\U^{heq})}$, the following is a pullback square.
    % https://q.uiver.app/#q=WzAsNCxbMCwwLCJGKFIpIl0sWzAsMSwiXFx0ZXh0dXB7Z3JhcGh9KGYpIl0sWzEsMCwiRihYXFx0aW1lcyBZKSJdLFsxLDEsIkYoWC9EXFx0aW1lcyBZL0UpIl0sWzAsMSwiIiwyLHsic3R5bGUiOnsiaGVhZCI6eyJuYW1lIjoiZXBpIn19fV0sWzIsMywiIiwwLHsic3R5bGUiOnsiaGVhZCI6eyJuYW1lIjoiZXBpIn19fV0sWzAsMiwiIiwwLHsic3R5bGUiOnsidGFpbCI6eyJuYW1lIjoiaG9vayIsInNpZGUiOiJ0b3AifX19XSxbMSwzLCIiLDAseyJzdHlsZSI6eyJ0YWlsIjp7Im5hbWUiOiJob29rIiwic2lkZSI6InRvcCJ9fX1dLFswLDMsIiIsMix7InN0eWxlIjp7Im5hbWUiOiJjb3JuZXIifX1dXQ==
    \[\begin{tikzcd}
    	{F(R)} & {F(X\times Y)} \\
    	{\textup{graph}(f)} & {F(X/D\times Y/E)}
    	\arrow[hook, from=1-1, to=1-2]
    	\arrow[two heads, from=1-1, to=2-1]
    	\arrow["\lrcorner"{anchor=center, pos=0.125}, draw=none, from=1-1, to=2-2]
    	\arrow[two heads, from=1-2, to=2-2]
    	\arrow[hook, from=2-1, to=2-2]
    \end{tikzcd}\]
    In $(\Pro(\defT(T)))_{ex/reg}$, the morphism $\tilde{f}$ being the morphism induced by $R$ means that we have the following diagram.
    % https://q.uiver.app/#q=WzAsNCxbMCwwLCJSIl0sWzAsMSwiXFx0ZXh0dXB7Z3JhcGh9KFxcdGlsZGV7Zn0pIl0sWzEsMCwiWFxcdGltZXMgWSJdLFsxLDEsIlgvRFxcdGltZXMgWS9FIl0sWzAsMSwiIiwyLHsic3R5bGUiOnsiaGVhZCI6eyJuYW1lIjoiZXBpIn19fV0sWzIsMywiIiwwLHsic3R5bGUiOnsiaGVhZCI6eyJuYW1lIjoiZXBpIn19fV0sWzAsMiwiIiwwLHsic3R5bGUiOnsidGFpbCI6eyJuYW1lIjoiaG9vayIsInNpZGUiOiJ0b3AifX19XSxbMSwzLCIiLDIseyJzdHlsZSI6eyJ0YWlsIjp7Im5hbWUiOiJob29rIiwic2lkZSI6InRvcCJ9fX1dLFswLDMsIiIsMix7InN0eWxlIjp7Im5hbWUiOiJjb3JuZXIifX1dXQ==
    \[\begin{tikzcd}
        R & {X\times Y} \\
        {\textup{graph}(\tilde{f})} & {X/D\times Y/E}
        \arrow[hook, from=1-1, to=1-2]
        \arrow[two heads, from=1-1, to=2-1]
        \arrow["\lrcorner"{anchor=center, pos=0.125}, draw=none, from=1-1, to=2-2]
        \arrow[two heads, from=1-2, to=2-2]
        \arrow[hook, from=2-1, to=2-2]
    \end{tikzcd}\]
    Hence, we have 
    \[\textup{graph}(F(\tilde{f})) \cong F(\textup{graph}(\tilde{f})) \cong \textup{graph}(f), \]
    where the first isomorphism comes from the regularity of the functor $F$, and the second isomorphism comes from the uniqueness of the regular-epimorphism-monomorphism factorisation in ${\tpdefT(\U^{heq})}$. Then, $F(\tilde{f}) = f$ by the universality of taking the graph of a morphism.

\smallskip

    Faithfulness: \quad Let $f,g\in\Hom_{(\Pro(\defT(T)))_{ex/reg}}(X/D,Y/E)$, so there are relations $R,S$ from $X$ to $Y$ in $\Pro(\defT(T))$ such that the following two squares are pullback squares.
    % https://q.uiver.app/#q=WzAsOCxbMCwwLCJSIl0sWzAsMSwiXFx0ZXh0dXB7Z3JhcGh9KGYpIl0sWzEsMCwiWFxcdGltZXMgWSJdLFsxLDEsIlgvRFxcdGltZXMgWS9FIl0sWzMsMCwiUyJdLFs0LDAsIlhcXHRpbWVzIFkiXSxbMywxLCJcXHRleHR1cHtncmFwaH0oZykiXSxbNCwxLCJYL0RcXHRpbWVzIFkvRSJdLFswLDEsIiIsMix7InN0eWxlIjp7ImhlYWQiOnsibmFtZSI6ImVwaSJ9fX1dLFsyLDMsIiIsMCx7InN0eWxlIjp7ImhlYWQiOnsibmFtZSI6ImVwaSJ9fX1dLFswLDIsIiIsMCx7InN0eWxlIjp7InRhaWwiOnsibmFtZSI6Imhvb2siLCJzaWRlIjoidG9wIn19fV0sWzEsMywiIiwyLHsic3R5bGUiOnsidGFpbCI6eyJuYW1lIjoiaG9vayIsInNpZGUiOiJ0b3AifX19XSxbMCwzLCIiLDIseyJzdHlsZSI6eyJuYW1lIjoiY29ybmVyIn19XSxbNSw3LCIiLDIseyJzdHlsZSI6eyJoZWFkIjp7Im5hbWUiOiJlcGkifX19XSxbNCw2LCIiLDIseyJzdHlsZSI6eyJoZWFkIjp7Im5hbWUiOiJlcGkifX19XSxbNiw3LCIiLDEseyJzdHlsZSI6eyJ0YWlsIjp7Im5hbWUiOiJob29rIiwic2lkZSI6InRvcCJ9fX1dLFs0LDUsIiIsMSx7InN0eWxlIjp7InRhaWwiOnsibmFtZSI6Imhvb2siLCJzaWRlIjoidG9wIn19fV0sWzQsNywiIiwxLHsic3R5bGUiOnsibmFtZSI6ImNvcm5lciJ9fV1d
    \[\begin{tikzcd}
        R & {X\times Y} && S & {X\times Y} \\
        {\textup{graph}(f)} & {X/D\times Y/E} && {\textup{graph}(g)} & {X/D\times Y/E}
        \arrow[hook, from=1-1, to=1-2]
        \arrow[two heads, from=1-1, to=2-1]
        \arrow["\lrcorner"{anchor=center, pos=0.125}, draw=none, from=1-1, to=2-2]
        \arrow[two heads, from=1-2, to=2-2]
        \arrow[hook, from=1-4, to=1-5]
        \arrow[two heads, from=1-4, to=2-4]
        \arrow["\lrcorner"{anchor=center, pos=0.125}, draw=none, from=1-4, to=2-5]
        \arrow[two heads, from=1-5, to=2-5]
        \arrow[hook, from=2-1, to=2-2]
        \arrow[hook, from=2-4, to=2-5]
    \end{tikzcd}\]
    Suppose $F(f)=F(g)$. Then in $\mathcal{U}$, for any $a,a'$ realising $X$ and any $b,b'$ realising $Y$, such that $R(a,b)$ and $S(a',b')$, if $D(a,a')$, then $E(b,b')$. Now, the relations $R,S$ satisfy Definition \ref{ex/reg carboni}. In particular, they are invariant under both equivalence relations $D,E$, and are total and single-valued over their quotients. 
    
    Suppose $R(a,b)$. As $S$ is total, $S(a,b')$ holds for some $b'$. By the above, $E(b,b')$ holds. $S$ is invariant under $E$, so $S(a,b)$ holds. The converse also holds by symmetry. Hence, $R=S$, so $f=g$.
\end{proof}

\begin{remark}
    In the positive setting, the functor $F$ still exists and is essentially surjective and faithful. However, as Example~\ref{positive not work ex}, it might not be full. The use of Lemma~\ref{functions of Pro} and Proposition~\ref{isomorphism is calculated in monster} in the proof for fullness prevents the proof from transferring to the positive setting.
\end{remark}

\medskip

\subsection{A negative result}
One might hope that there is some free construction on $T$, giving us a first-order (or even positive) theory that eliminates hyperimaginaries, similar to how $T^{eq}$ is the free first-order theory on $T$ that eliminates imaginaries. Categorically, the hope is that there is some subcategory $\C$ of $(\Pro(\defT(T)))_{ex/reg}$, such that $\Pro(\C) \simeq (\Pro(\defT(T)))_{ex/reg}$. If such a subcategory exists, then we can think of its internal theory as the free construction on $T$ that eliminates hyperimaginaries. Alas, this is not true. We show such a subcategory does not exist for Example~\ref{finset}, that is, there is no category, whose pro-completion is \textbf{CHaus}. Now, the dual of this statement---$\textbf{CHaus}^{op}$ is not the ind-completion of any category---is well-known. For example, it can be deduced from \cite{LIEBERMAN2023107245}, where they showed that $\textbf{CHaus}^{op}$ is not finitely concrete. We give a more direct proof here.

By Remark~\ref{cocompact}, if such a subcategory of \textbf{CHaus} exists, then it has to be the full subcategory consisting of the cocompact objects. We have the following characterisation of cocompact objects in \textbf{CHaus}. This is folklore, but it seems that the proof is not written down anywhere.

\begin{prop}
    The cocompact objects of \textup{\textbf{CHaus}} are precisely the finite discrete spaces.
\end{prop}

\begin{proof}
    $(\Rightarrow)$ \quad Suppose $X\in\textbf{CHaus}$ is cocompact. Let $X_d$ be the space with the same underlying set as $X$ and the discrete topology. Let $\beta X_d$ be the Stone-\v{C}ech compactification of $X_d$. Then, there exists a unique continuous function $f:\beta X_d\rightarrow X$ extending the obvious morphism $\iota:X_d\rightarrow X$. Now, $\beta X_d$ is a Stone space, so it is a cofiltered limit of finite discrete spaces, say $\beta X_d\cong \lim F_i$, where each $F_i$ is finite. Since $X$ is cocompact, $f$ factors through some $F_i$. Hence, we obtain the following diagram.
    % https://q.uiver.app/#q=WzAsNCxbMCwwLCJYX2QiXSxbMSwwLCJYIl0sWzAsMSwiXFxiZXRhIFhfZCJdLFsxLDEsIkZfaSJdLFswLDEsIlxcaW90YSJdLFswLDJdLFsyLDEsImYiLDFdLFszLDEsImZfaSIsMl0sWzIsMywiXFxwaV9pIiwyXV0=
    \[\begin{tikzcd}
    	{X_d} & X \\
    	{\beta X_d} & {F_i}
    	\arrow["\iota", from=1-1, to=1-2]
    	\arrow[from=1-1, to=2-1]
    	\arrow["f"{description}, from=2-1, to=1-2]
    	\arrow["{\pi_i}"', from=2-1, to=2-2]
    	\arrow["{f_i}"', from=2-2, to=1-2]
    \end{tikzcd}\]
    Since $\iota$ is a surjection, $f_i$ is also surjective. Since $F_i$ is finite, $X$ also has to be finite. 

    \smallskip

    $(\Leftarrow)$ \quad Let $F$ be a finite set with the discrete topology. Let $(X_i)_{i\in I}$ be a cofiltered system in \textbf{CHaus}, let $X$ denote the limit, let $p_i:X\rightarrow X_i$ denote the projection, and let $f:X \rightarrow F$ be a continuous function. Since $F$ is finite discrete, we can write $X$ as the disjoint union of the fibres $f^{-1}(a)$ for $a\in F$. Each fibre is clopen.
    
    Now, $X$ has the limit topology, so has basis given by $\{p_i^{-1}(V) \ | \ i\in I, V\subseteq X_i \textup{ open}\}$. For each $x\in X$, we choose some basic open $p_{i_x}^{-1}(V_x)\ni x$, such that $p_{i_x}^{-1}(V_x) \subseteq f^{-1}(f(x))$. This forms a cover of $X$. Since $X$ is compact, there is a finite subcover $\{p_{i_1}^{-1}(V_1),\ldots,p_{i_n}^{-1}(V_n)\}$.
    
    Since $I$ is cofiltered, there exists $i\in I$ with transition maps $p_{ii_r}:X_i\rightarrow X_{i_r}$ for $r=1,\ldots,n$, such that $p_{ii_r}p_i=p_{i_r}$. For each $r$, let $W_r = p_{ii_r}^{-1}(V_r) \subseteq X_i$. Then $p_i^{-1}(W_r) = p_{i_r}^{-1}(V_r)$, so $\{p_i^{-1}(W_r)\}_{r=1}^n$ is a cover of $X$. Moreover, since we have chosen each basic open to be contained in a single fibre of $f$, there exists a unique continuous map $h:\bigcup_{r=1}^n W_r \rightarrow F$ extending $f$ along $p_i$. 

    Note that $p_i(X)\subseteq W:= \bigcup W_r \subseteq X_i$. We show that there is some $j$ with a transition map $p_{ji}:X_j\rightarrow X_i$, whose image is contained in $W$. Otherwise, for all $j\in I$ with a morphism $j\rightarrow i$, the closed sets $p_{ji}(X_j)-W$ are non-empty and form a cofiltered system of compact Hausdorff spaces. Therefore, the limit, which is a subset of $X$, is non-empty, contradicting the fact that $p_i(X)\subseteq W$. Hence, fix $j$ such that $p_{ji}(X_j)\subseteq W$ and write $p_{ji}^W:X_j\rightarrow W$ for the same map with codomain restricted to $W$. Then, $f$ factors through $X_j$ via the map $h p^W_{ji}:X_j \rightarrow F$.
% https://q.uiver.app/#q=WzAsOCxbMiwwLCJGIl0sWzAsMCwiWCJdLFswLDEsIlhfaiJdLFsxLDIsIlciXSxbMCwyLCJYX2kiXSxbMCwzLCJYX3tpX2t9Il0sWzIsMywiVl9rIl0sWzIsMiwiV19rIl0sWzEsMCwiZiJdLFsxLDIsInBfaiIsMl0sWzIsMywicF97aml9XlciXSxbMiw0LCJwX3tqaX0iLDJdLFszLDAsImgiLDJdLFs0LDUsInBfe2lpX2t9IiwyXSxbMyw0LCIiLDEseyJzdHlsZSI6eyJ0YWlsIjp7Im5hbWUiOiJob29rIiwic2lkZSI6InRvcCJ9fX1dLFs2LDUsIiIsMCx7InN0eWxlIjp7InRhaWwiOnsibmFtZSI6Imhvb2siLCJzaWRlIjoidG9wIn19fV0sWzcsNl0sWzcsNSwiIiwwLHsic3R5bGUiOnsibmFtZSI6ImNvcm5lciJ9fV0sWzcsMywiIiwyLHsic3R5bGUiOnsidGFpbCI6eyJuYW1lIjoiaG9vayIsInNpZGUiOiJ0b3AifX19XV0=
\[\begin{tikzcd}
	X && F \\
	{X_j} \\
	{X_i} & W & {W_r} \\
	{X_{i_r}} && {V_r}
	\arrow["f", from=1-1, to=1-3]
	\arrow["{p_j}"', from=1-1, to=2-1]
	\arrow["{p_{ji}}"', from=2-1, to=3-1]
	\arrow["{p_{ji}^W}", from=2-1, to=3-2]
	\arrow["{p_{ii_r}}"', from=3-1, to=4-1]
	\arrow["h"', from=3-2, to=1-3]
	\arrow[hook, from=3-2, to=3-1]
	\arrow[hook, from=3-3, to=3-2]
	\arrow["\lrcorner"{anchor=center, pos=0.125, rotate=-90}, draw=none, from=3-3, to=4-1]
	\arrow[from=3-3, to=4-3]
	\arrow[hook, from=4-3, to=4-1]
\end{tikzcd}\]
\end{proof}

The proposition shows that we have $\textup{Cocompact}(\textbf{CHaus}) \simeq \textbf{FinSet}$, whose pro-completion is \textbf{Stone}. In particular, there is no subcategory of \textbf{CHaus}, whose pro-completion is \textbf{CHaus} itself. 

\medskip



We conclude by comparing a few categories obtained by applying the pro-completion and the exact completion to $\defT(T)$ in different ways. The following diagram is obtained by inducing the dashed arrows in the labelled order using the universal properties of the pro-completion and the exact completion.
% https://q.uiver.app/#q=WzAsOCxbMSwwLCJcXGRlZlQoVCkiXSxbMSwxLCJcXFBybyhcXGRlZlQoVCkpIl0sWzEsMiwiKFxcUHJvKFxcZGVmVChUKSkpX3tleC9yZWd9Il0sWzMsMCwiXFxkZWZUKFQpX3tleC9yZWd9Il0sWzQsMCwiXFxkZWZUKFRee2VxfSkiXSxbMywxLCJcXFBybyAoXFxkZWZUKFQpX3tleC9yZWd9KSJdLFszLDIsIihcXFBybyAoXFxkZWZUKFQpX3tleC9yZWd9KSlfe2V4L3JlZ30iXSxbMCwyLCJcXHRwZGVmVChcXFVee2hlcX0pIl0sWzAsMV0sWzEsMl0sWzAsM10sWzMsNCwiXFxzaW1lcSIsMSx7InN0eWxlIjp7ImJvZHkiOnsibmFtZSI6Im5vbmUifSwiaGVhZCI6eyJuYW1lIjoibm9uZSJ9fX1dLFszLDVdLFs1LDZdLFsxLDUsIjEiLDEseyJsYWJlbF9wb3NpdGlvbiI6MzAsInN0eWxlIjp7ImJvZHkiOnsibmFtZSI6ImRhc2hlZCJ9fX1dLFs1LDIsIjQiLDEseyJzdHlsZSI6eyJib2R5Ijp7Im5hbWUiOiJkYXNoZWQifX19XSxbNywyLCJcXHNpbWVxIiwxLHsic3R5bGUiOnsiYm9keSI6eyJuYW1lIjoibm9uZSJ9LCJoZWFkIjp7Im5hbWUiOiJub25lIn19fV0sWzMsMiwiMyIsMSx7ImxhYmVsX3Bvc2l0aW9uIjoyMCwic3R5bGUiOnsiYm9keSI6eyJuYW1lIjoiZGFzaGVkIn19fV0sWzIsNiwiMiIsMSx7Im9mZnNldCI6LTMsInN0eWxlIjp7ImJvZHkiOnsibmFtZSI6ImRhc2hlZCJ9fX1dLFs2LDIsIjUiLDEseyJvZmZzZXQiOi0zLCJzdHlsZSI6eyJib2R5Ijp7Im5hbWUiOiJkYXNoZWQifX19XSxbMiw2LCJcXHNpbSIsMSx7InN0eWxlIjp7ImJvZHkiOnsibmFtZSI6Im5vbmUifSwiaGVhZCI6eyJuYW1lIjoibm9uZSJ9fX1dXQ==
\[\begin{tikzcd}
	& {\defT(T)} && {\defT(T)_{ex/reg}} & {\defT(T^{eq})} \\
	& {\Pro(\defT(T))} && {\Pro (\defT(T)_{ex/reg})} \\
	{\tpdefT(\U^{heq})} & {(\Pro(\defT(T)))_{ex/reg}} && {(\Pro (\defT(T)_{ex/reg}))_{ex/reg}}
	\arrow[from=1-2, to=1-4]
	\arrow[from=1-2, to=2-2]
	\arrow["\simeq"{description}, draw=none, from=1-4, to=1-5]
	\arrow[from=1-4, to=2-4]
	\arrow["3"{description, pos=0.2}, dashed, from=1-4, to=3-2]
	\arrow["1"{description, pos=0.3}, dashed, from=2-2, to=2-4]
	\arrow[from=2-2, to=3-2]
	\arrow["4"{description}, dashed, from=2-4, to=3-2]
	\arrow[from=2-4, to=3-4]
	\arrow["\simeq"{description}, draw=none, from=3-1, to=3-2]
	\arrow["2"{description}, shift left=3, dashed, from=3-2, to=3-4]
	\arrow["\sim"{description}, draw=none, from=3-2, to=3-4]
	\arrow["5"{description}, shift left=3, dashed, from=3-4, to=3-2]
\end{tikzcd}\]

The bottom row is always an equivalence. If $T$ eliminates hyperimaginaries, then the bottom right three categories are equivalent. If $T$ eliminates imaginaries, then the three rows are equivalences. If $T$ eliminates both imaginaries and hyperimaginaries, then the bottom four categories are equivalent. Hence, we obtain our main theorem.

\begin{thm}
    
    Let $T$ be a consistent complete first-order theory and let $\U$ be a monster model.
    \begin{enumerate}
    \item $T$ eliminates imaginaries if and only if $\defT(T)$ is exact.

    \item $T$ eliminates hyperimaginaries if and only if $\Pro(\defT(T)_{ex/reg})$ is exact.
        
    \item $T$ eliminates imaginaries and hyperimaginaries if and only if $\Pro(\defT(T))$ is exact.

    \item We have the following equivalences of categories: 
        \begin{align*}
            \defT(T^{eq}) &\simeq \defT(T)_{ex/reg}, \\
            \tpdefT(\U^{heq}) &\simeq (\Pro(\defT(T)))_{ex/reg}.
        \end{align*}
    \end{enumerate}


\end{thm}



